% !TEX root = ../notes.tex

\documentclass[../notes.tex]{subfiles}

\begin{document}

\section{April 8}
Here we go.

\subsection{Purity}
Another property of Brauer groups is purity. For motivation, note that any normal integral Noetherian scheme $X$ admits an exact sequence
\[1\to\OO(X)^\times\to K(X)^\times\to\bigoplus_{x\in X^{(1)}}\ZZ,\]
where $X^{(1)}$ is the set of codimension $1$ points. Namely, we are asserting that a meromorphic function $f\in K(X)^\times$ is an invertible regular function on $X$ if and only if it has vanishing valuation at all codimension $1$ points. The analog of this for Brauer groups is as follows.
\begin{theorem} \label{thm:harthog-for-brauer}
	Fix a regular integral Noetherian scheme $X$. Then there is an exact sequence
	\[0\to\op{Br}X\to\op{Br}K(X)\to\bigoplus_{x\in X^{(1)}}\mathrm H^1(k(x);\QQ/\ZZ),\]
	where one implicitly removes all the $p$-primary parts if any of the $k(x)$ are imperfect of characteristic $p$.
\end{theorem}
\begin{remark}
	We will not prove this or even define the maps on the right, but we will say that this recovers the left-exactness of \Cref{ex:cft-exact-seq} when $k$ is a function field.
\end{remark}
\begin{corollary} \label{cor:extend-brauer}
	Fix a regular integral Noetherian variety $X$ over a field of characteristic zero. Given a Zariski open cover $\{U_i\}$ of a regular integral Noetherian scheme $X$,
	\[\op{Br}X=\bigcap_i\op{Br}U_i,\]
	where the intersection takes place in $\op{Br}K(X)$.
\end{corollary}
\begin{proof}
	To begin, use \Cref{thm:harthog-for-brauer} to present $\op{Br}X$ as elements of $\op{Br}K(X)$ which vanish in the various groups $\mathrm H^1(k(x);\QQ/\ZZ)$. Then the right-hand side is also presented as elements of $\op{Br}X$ which vanish in the various groups $\mathrm H^1(k(x);\QQ/\ZZ)$ where $x\in U_i$, so we recover the result by looping over all $U_i$s.
\end{proof}
One can remove the characteristic zero assumption, but it is harder. The following was conjectured by Grothendieck in 1968 but only proven in 2019 with perfectoid methods.
\begin{theorem}[\v{C}esnavi\v{c}ius] \label{thm:brauer-purity}
	Fix a regular integral Noetherian scheme $X$. For a closed subscheme $Z\subseteq X$ of codimension at least $2$, set $U\coloneqq X\setminus Z$. Then the natural map $\op{Br}X\to\op{Br}U$ is an isomorphism.
\end{theorem}
\begin{remark}
	This follows from \Cref{thm:harthog-for-brauer}, except for the $p$-primary parts.
\end{remark}
\begin{corollary} \label{cor:birational-vars}
	Fix nice birational varieties $X$ and $Y$. Then $\op{Br}X\cong\op{Br}Y$.
\end{corollary}
\begin{proof}
	Choose our birational map $f\colon X\to Y$. Then $f$ induces an isomorphism $K(X)\to K(Y)$, so
	\[\op{Br}K(X)=\op{Br}K(Y).\]
	Now, $f$ is defined on some open subset $U\subseteq X$, and we may even assume that the codimension of $X\setminus U$ is at least $2$. Thus, there is a well-defined map $\op{Br}Y\to\op{Br}U$, which is then isomorphic to $\op{Br}X$ (by \Cref{thm:brauer-purity}!), which then embeds back into $\op{Br}K(X)=\op{Br}K(Y)$. Thus, $\op{Br}Y\subseteq\op{Br}X$ as subgroups of the Brauer group of the function field. The opposite inclusion holds by symmetry, so we are done.
\end{proof}

\subsection{Examples of Brauer Groups of Schemes} \label{subsec:brauer-computations}
Let's do an example.
\begin{example} \label{ex:br-p1-field}
	For a field $k$, we claim that $\op{Br}\PP^1_k=\op{Br}k$. Indeed, $\op{Pic}\PP^1_k=\PP^1_{k^{\mathrm{sep}}}=\ZZ$, so \Cref{cor:brauer-les} gives us the exact sequence
	\[0\to\ZZ\to\ZZ^G\to\op{Br}k\to\op{Br}_1\PP^1_k\to\mathrm H^1(G;\ZZ).\]
	Now, because $G$ acts trivially on $\ZZ$, we see $\mathrm H^1(G;\ZZ)=\op{Hom}_{\mathrm{cts}}(G,\ZZ)$; because $G$ is compact, any map $G\to\ZZ$ has compact image, but the only compact subgroup of $\ZZ$ is trivial, so we conclude $\mathrm H^1(G;\ZZ)=0$. Thus, we see that we have an isomorphism $\op{Br}k\to\op{Br}_1\PP^1_k$. However, $\op{Br}_1\PP^1_k$ is the kernel of the map $\op{Br}\PP^1_k\to\op{Br}\PP^1_{k^{\mathrm{sep}}}$, and the latter vanishes by \Cref{rem:br-k-sep-curve}.
\end{example}
\begin{remark}
	The group $\op{Br}X$ can be much more complicated for higher genus curves or higher-dimensional varieties because $\op{Pic}X$ is more complicated.
\end{remark}
We can upgrade \Cref{ex:br-p1-field} as follows.
\begin{proposition} \label{prop:br-relative-p1}
	Fix a regular integral Noetherian scheme $B$. Then the natural map $\op{Br}B\to\op{Br}\PP^1_B$ is an isomorphism.
\end{proposition}
\begin{proof}
	There is a pullback square
	% https://q.uiver.app/#q=WzAsNCxbMCwwLCJcXFBQXjFfe0soQil9Il0sWzEsMCwiXFxQUF4xX0IiXSxbMSwxLCJCIl0sWzAsMSwiXFxTcGVjIEsoQikiXSxbMywyXSxbMCwzXSxbMCwxXSxbMSwyXV0=&macro_url=https%3A%2F%2Fraw.githubusercontent.com%2FdFoiler%2Fnotes%2Fmaster%2Fnir.tex
	\[\begin{tikzcd}[cramped]
		{\PP^1_{K(B)}} & {\PP^1_B} \\
		{\Spec K(B)} & B
		\arrow[from=1-1, to=1-2]
		\arrow[from=1-1, to=2-1]
		\arrow[from=1-2, to=2-2]
		\arrow[from=2-1, to=2-2]
	\end{tikzcd}\]
	which on Brauer groups produces the following commutative diagram.
	% https://q.uiver.app/#q=WzAsNSxbMCwwLCJcXG9we0JyfUIiXSxbMSwwLCJcXG9we0JyfUsoQikiXSxbMCwxLCJcXG9we0JyfVxcUFBeMV9CIl0sWzEsMSwiXFxvcHtCcn1cXFBQXjFfe0soQil9Il0sWzIsMSwiXFxvcHtCcn1LXFxsZWZ0KFxcUFBeMV97SyhCKX1cXHJpZ2h0KSJdLFswLDFdLFsxLDNdLFsyLDNdLFszLDRdLFswLDJdXQ==&macro_url=https%3A%2F%2Fraw.githubusercontent.com%2FdFoiler%2Fnotes%2Fmaster%2Fnir.tex
	\[\begin{tikzcd}[cramped]
		{\op{Br}B} & {\op{Br}K(B)} & \\
		{\op{Br}\PP^1_B} & {\op{Br}\PP^1_{K(B)}} & {\op{Br}K\left(\PP^1_{K(B)}\right)}
		\arrow[from=1-1, to=1-2]
		\arrow[from=1-1, to=2-1]
		\arrow[from=1-2, to=2-2]
		\arrow[from=2-1, to=2-2]
		\arrow[from=2-2, to=2-3]
	\end{tikzcd}\]
	Here, the top arrow and the bottom composite are both embeddings, so we conclude that all horizontal maps are embeddings. Now, the right vertical map is an isomorphism by \Cref{ex:br-p1-field}. The commutativity of the diagram then implies that $\op{Br}B$ must be embedded in $\op{Br}\PP^1_B$.
	
	Conversely, note that the projection $\PP^1_B\to B$ admits easy sections $s$, which then must give the inverse map $\op{Br}\PP^1_{K(B)}\to\op{Br}K(B)$. Thus, the inverse isomorphism of the right vertical arrow also fits into a commutative diagram
	% https://q.uiver.app/#q=WzAsNSxbMCwwLCJcXG9we0JyfUIiXSxbMSwwLCJcXG9we0JyfUsoQikiXSxbMCwxLCJcXG9we0JyfVxcUFBeMV9CIl0sWzEsMSwiXFxvcHtCcn1cXFBQXjFfe0soQil9Il0sWzIsMSwiXFxvcHtCcn1LXFxsZWZ0KFxcUFBeMV97SyhCKX1cXHJpZ2h0KSJdLFswLDFdLFszLDEsInMiXSxbMiwzXSxbMyw0XSxbMiwwLCJzIl1d&macro_url=https%3A%2F%2Fraw.githubusercontent.com%2FdFoiler%2Fnotes%2Fmaster%2Fnir.tex
	\[\begin{tikzcd}[cramped]
		{\op{Br}B} & {\op{Br}K(B)} & \\
		{\op{Br}\PP^1_B} & {\op{Br}\PP^1_{K(B)}} & {\op{Br}K\left(\PP^1_{K(B)}\right)}
		\arrow[from=1-1, to=1-2]
		\arrow["s", from=2-1, to=1-1]
		\arrow[from=2-1, to=2-2]
		\arrow["s", from=2-2, to=1-2]
		\arrow[from=2-2, to=2-3]
	\end{tikzcd}\]
	so it follows that $\op{Br}\PP^1_B$ embeds into $\PP^1_{K(B)}$, as required!
\end{proof}
\begin{example} \label{ex:br-pn}
	It follows from \Cref{prop:br-relative-p1} that $\op{Br}\left(\PP^1_k\times X\right)\cong\op{Br}X$ for a nice $k$-variety $X$. For example, $\PP^n_k$ is birational to iterated products of $\PP^1_k$, so \Cref{cor:birational-vars} allows us to conclude that $\op{Br}\PP^n_k=\op{Br}k$.
\end{example}
Here is another example.
\begin{definition}[quadric]
	A \textit{quadric hypersurface} $X$ over a field $k$ is a degree $2$ hypersurface of some projective space $\PP^n_k$.
\end{definition}
\begin{proposition} \label{prop:br-quadric-hypersurface}
	Fix a smooth quadric hypersurface $X$ over a field $k$. Then the natural map $\op{Br}k\to\op{Br}X$ is surjective.
\end{proposition}
\begin{proof}
	The main point is to use the exact sequence
	\[\op{Br}k\to\op{Br}_1X\to\mathrm H^1(G;\op{Pic}X_{k^{\mathrm{sep}}})\]
	from \Cref{cor:brauer-les}.
	\begin{itemize}
		\item We claim $\op{Br}_1X=\op{Br}X$. One can show (by putting the quadratic equation in a standard form) that $X_{k^{\mathrm{sep}}}$ is birational to $\PP^{n-1}_{k^{\mathrm{sep}}}$. Thus, by \Cref{cor:birational-vars} and \Cref{ex:br-pn}, we see that $\op{Br}X_{k^{\mathrm{sep}}}=0$, so $\op{Br}_1X=\op{Br}X$.
		\item We claim that $\mathrm H^1(G;\op{Pic}X_{k^{\mathrm{sep}}})=0$. One can compute (it's \cite[Exercise~II.6.5]{hartshorne}) that
		\[\op{Pic}X_{k^{\mathrm{sep}}}=\begin{cases}
			\ZZ^2 & \text{if }\dim X=2, \\
			\ZZ & \text{if }\dim X=2.
		\end{cases}\]
		In the case where the group is $\ZZ$, the Galois group should act trivially: the Galois group needs to preserve the fact that a line bundle is ample, so it needs to preserve the positives, so it needs to act by identity. In this case, one concludes that $\mathrm H^1(G;\ZZ)=0$ as in \Cref{ex:br-p1-field}.
		
		Even if the Picard group is $\ZZ^2$, one can compute the Galois cohomology must vanish, so we are done. Indeed, in this case, $X_{k^{\mathrm{sep}}}$ is isomorphic to $\PP^1_{k^{\mathrm{sep}}}\times\PP^1_{k^{\mathrm{sep}}}$. If $G$ acts trivially on $\ZZ^2$, then we argue as in the previous paragraph. Otherwise, $G$ interchanges the two generators, but $\mathrm H^1(G;\op{Pic}X_{k^{\mathrm{sep}}})=0$ in this case too.
		\qedhere
	\end{itemize}
\end{proof}
We can also do a relative version.
\begin{definition}[quadric bundle]
	A \textit{quadric bundle} is a morphism $\pi\colon X\to B$ is a flat morphism of regular integral $k$-varieties satisfying the following.
	\begin{itemize}
		\item Quadric: the generic fiber $X_{K(B)}$ is a smooth quadric.
		\item Every fiber of $\pi$ has an irreducible component of multiplicity $1$ that is geometrically irreducible.
	\end{itemize}
\end{definition}
\begin{remark}
	The second point basically prevents a fiber from looking like $X^2-2Y^2$ for a quadric bundle $\pi\colon X\to\Spec\ZZ$. The moral is that 
\end{remark}
\begin{proposition} \label{prop:br-quadric-bundle}
	Fix a quadric bundle $\pi\colon X\to B$ over a field $k$ of characteristic zero. Then the natural map $\op{Br}B\to\op{Br}X$ is surjective.
\end{proposition}
\begin{proof}
	We use \Cref{thm:harthog-for-brauer}, which gives us a morphism
	% https://q.uiver.app/#q=WzAsOCxbMCwwLCIwIl0sWzEsMCwiXFxvcHtCcn1CIl0sWzIsMCwiXFxvcHtCcn1LKEIpIl0sWzIsMSwiXFxvcHtCcn1LKFgpIl0sWzEsMSwiXFxvcHtCcn1YIl0sWzAsMSwiMCJdLFszLDAsIlxcZGlzcGxheXN0eWxlXFxiaWdvcGx1c197YlxcaW4gQl57KDEpfX1cXG1hdGhybSBIXjEoayhiKTtcXFFRL1xcWlopIl0sWzMsMSwiXFxkaXNwbGF5c3R5bGVcXGJpZ29wbHVzX3t4XFxpbiBYXnsoMSl9fVxcbWF0aHJtIEheMSh4KGIpO1xcUVEvXFxaWikiXSxbMiwzXSxbNSw0XSxbMCwxXSxbMSwyXSxbNCwzXSxbMSw0XSxbMiw2LCJcXG9we3Jlc30iXSxbMyw3LCJcXG9we3Jlc30iXV0=&macro_url=https%3A%2F%2Fraw.githubusercontent.com%2FdFoiler%2Fnotes%2Fmaster%2Fnir.tex
	\[\begin{tikzcd}[cramped]
		0 & {\op{Br}B} & {\op{Br}K(B)} & {\displaystyle\bigoplus_{b\in B^{(1)}}\mathrm H^1(k(b);\QQ/\ZZ)} \\
		0 & {\op{Br}X} & {\op{Br}K(X)} & {\displaystyle\bigoplus_{x\in X^{(1)}}\mathrm H^1(k(x);\QQ/\ZZ)}
		\arrow[from=1-1, to=1-2]
		\arrow[from=1-2, to=1-3]
		\arrow[from=1-2, to=2-2]
		\arrow["{\op{res}}", from=1-3, to=1-4]
		\arrow[from=1-3, to=2-3]
		\arrow[from=2-1, to=2-2]
		\arrow[from=2-2, to=2-3]
		\arrow["{\op{res}}", from=2-3, to=2-4]
	\end{tikzcd}\]
	of exact sequences. In fact, we can compare the analogous exact sequence for $\op{Br}X_{K(B)}$, which we see is then cut out by those codimension $1$ points $x\in X^{(1)}$ which surject onto $\Spec B$. Thus, by looping over the remaining points in $X^{(1)}$, we get a morphism
	% https://q.uiver.app/#q=WzAsOCxbMCwwLCIwIl0sWzEsMCwiXFxvcHtCcn1CIl0sWzIsMCwiXFxvcHtCcn1LKEIpIl0sWzIsMSwiXFxvcHtCcn1LKFgpIl0sWzEsMSwiXFxvcHtCcn1YIl0sWzAsMSwiMCJdLFszLDAsIlxcZGlzcGxheXN0eWxlXFxiaWdvcGx1c197YlxcaW4gQl57KDEpfX1cXG1hdGhybSBIXjEoayhiKTtcXFFRL1xcWlopIl0sWzMsMSwiXFxkaXNwbGF5c3R5bGVcXGJpZ29wbHVzX3tcXHN1YnN0YWNre2JcXGluIEJeeygxKX1cXFxcXFxwaSh4KT1ifX1cXG1hdGhybSBIXjEoayh4KTtcXFFRL1xcWlopIl0sWzIsM10sWzUsNF0sWzAsMV0sWzEsMl0sWzQsM10sWzEsNF0sWzIsNiwiXFxvcHtyZXN9Il0sWzMsNywiXFxvcHtyZXN9Il0sWzYsN11d&macro_url=https%3A%2F%2Fraw.githubusercontent.com%2FdFoiler%2Fnotes%2Fmaster%2Fnir.tex
	\[\begin{tikzcd}[cramped]
		0 & {\op{Br}B} & {\op{Br}K(B)} & {\displaystyle\bigoplus_{b\in B^{(1)}}\mathrm H^1(k(b);\QQ/\ZZ)} \\
		0 & {\op{Br}X} & {\op{Br}K(X)} & \begin{array}{c} \displaystyle\bigoplus_{\substack{b\in B^{(1)}\\\pi(x)=b}}\mathrm H^1(k(x);\QQ/\ZZ) \end{array}
		\arrow[from=1-1, to=1-2]
		\arrow[from=1-2, to=1-3]
		\arrow[from=1-2, to=2-2]
		\arrow["{\op{res}}", from=1-3, to=1-4]
		\arrow[from=1-3, to=2-3]
		\arrow[from=1-4, to=2-4]
		\arrow[from=2-1, to=2-2]
		\arrow[from=2-2, to=2-3]
		\arrow["{\op{res}}", from=2-3, to=2-4]
	\end{tikzcd}\]
	of exact sequences. By some diagram chasing (with the Four lemma), it is enough to show that the right vertical morphism is injective because we already know that the middle vertical morphism is surjective by \Cref{prop:br-quadric-hypersurface}.
	
	To this end, for a given $b\in B^{(1)}$ we choose a generic point $x$ of the hypersurface granted from the second assumption on our quadric bundle. Then $k(x)$ is an extension of $k(b)$ of some transcendence degree, so there is a surjection of Galois groups $\op{Gal}(k(x)^{\mathrm{sep}}/k(x))\to\op{Gal}(k(b)^{\mathrm{sep}}/k(b))$, which then induces an embedding on Galois cohomology. The result follows.
\end{proof}

\subsection{The Brauer--Manin Obstruction}
% It will be helpful to view $\op{Br}$ as a function we can evaluate.
\begin{definition}
	% Fix a morphism $\pi\colon X\to B$. For $\alpha\in\op{Br}X$, we define the \textit{evaluation map} $\op{ev}_\alpha\colon X(B)\to\op{Br}B$ by
	% \[\op{ev}_\alpha(x)\coloneqq x^*\alpha,\]
	% where we view $x\in X(k)$ as a morphism $\Spec k\to\Spec X$.
	Fix a $k$-variety $X$, and choose some $\alpha\in\op{Br}X$. For any $k$-algebra $L$, we define the map $\op{ev}_\alpha\colon X(L)\to\op{Br}L$ by sending $x\in X(L)$ to $x^*\alpha$, where we view $x$ as a morphism $\Spec L\to X$.
\end{definition}
% \begin{remark}
% 	In most applications, we will take $B=\Spec k$ so that $\op{ev}_\alpha$ is basically ``evaluating'' $\alpha$ on the $k$-points of $X$.
% \end{remark}
\begin{proposition} \label{prop:evaluation-vanishes-ae}
	Fix a nice $k$-variety $X$ over a global field $k$, and choose some $\alpha\in\op{Br}X$. Then the evaluation map $\op{ev}_\alpha\colon X(k_v)\to\op{Br}k_v$ vanishes for all but finitely many $v$.
\end{proposition}
\begin{proof}
	Choose a set $S$ of finite places large enough so that $X$ spreads out to some nice $k$-variety $\mf X$. By expanding $S$, we may assume that $\alpha$ extends to some $\widetilde\alpha\in\op{Br}\mf X$: indeed, we may simply add to $S$ the denominators needed to be able to define the Azumaya algebra corresponding to $\alpha$.

	We now claim that $\op{ev}_\alpha$ vanishes for $v\notin S$. Indeed, note that the diagram
	% https://q.uiver.app/#q=WzAsNCxbMCwwLCJcXG1mIFgoXFxPT192KSJdLFsxLDAsIlxcb3B7QnJ9XFxPT192Il0sWzAsMSwiWChrX3YpIl0sWzEsMSwiXFxvcHtCcn1rX3YiXSxbMCwxLCJcXG9we2V2fV97XFx3aWRldGlsZGVcXGFscGhhfSJdLFsxLDNdLFsyLDMsIlxcb3B7ZXZ9X1xcYWxwaGEiXSxbMCwyXV0=&macro_url=https%3A%2F%2Fraw.githubusercontent.com%2FdFoiler%2Fnotes%2Fmaster%2Fnir.tex
	\[\begin{tikzcd}[cramped]
		{\mf X(\OO_v)} & {\op{Br}\OO_v} \\
		{X(k_v)} & {\op{Br}k_v}
		\arrow["{\op{ev}_{\widetilde\alpha}}", from=1-1, to=1-2]
		\arrow[from=1-1, to=2-1]
		\arrow[from=1-2, to=2-2]
		\arrow["{\op{ev}_\alpha}", from=2-1, to=2-2]
	\end{tikzcd}\]
	commutes, so it is enough to recall that $\mf X(\OO_v)=X(k_v)$ by properness, and $\op{Br}\OO_v=0$ by \Cref{prop:br-dvr}.
\end{proof}
We are now ready to describe the Brauer--Manin obstruction. We continue with a nice $k$-variety $X$ and $\alpha\in\op{Br}X$. By the valuative criterion for properness, $X(\AA_k)=\prod_vX(k_v)$. Then \Cref{prop:evaluation-vanishes-ae} tells us that $\op{ev}_\alpha\colon X(\AA_k)\to\prod_v\op{Br}k_v$ actually lands in $\bigoplus_v\op{Br}k_v$, so the diagram
% https://q.uiver.app/#q=WzAsNyxbMSwwLCJYKGspIl0sWzIsMCwiWChcXEFBX2spIl0sWzIsMSwiXFxkaXNwbGF5c3R5bGVcXGJpZ29wbHVzX3ZcXG9we0JyfWtfdiJdLFsxLDEsIlxcb3B7QnJ9ayJdLFswLDEsIjAiXSxbMywxLCJcXFFRL1xcWloiXSxbNCwxLCIwIl0sWzEsMiwiXFxvcHtldn1fXFxhbHBoYSJdLFswLDFdLFswLDMsIlxcb3B7ZXZ9X1xcYWxwaGEiLDJdLFszLDJdLFs0LDNdLFsyLDVdLFs1LDZdXQ==&macro_url=https%3A%2F%2Fraw.githubusercontent.com%2FdFoiler%2Fnotes%2Fmaster%2Fnir.tex
\[\begin{tikzcd}[cramped]
	& {X(k)} & {X(\AA_k)} && \\
	0 & {\op{Br}k} & {\displaystyle\bigoplus_v\op{Br}k_v} & {\QQ/\ZZ} & 0
	\arrow[from=1-2, to=1-3]
	\arrow["{\op{ev}_\alpha}"', from=1-2, to=2-2]
	\arrow["{\op{ev}_\alpha}", from=1-3, to=2-3]
	\arrow[from=2-1, to=2-2]
	\arrow[from=2-2, to=2-3]
	\arrow[from=2-3, to=2-4]
	\arrow[from=2-4, to=2-5]
\end{tikzcd}\]
commutes, where the bottom row comes from global class field theory. Notably, a rational point $x\in X(k)$ not only gives to a point in $X(\AA_k)$, but we see that class must also vanish in $\QQ/\ZZ$ by the commutative diagram.
\begin{definition}[Brauer set]
	Fix a nice variety $X$ over a global field $k$. For each $\alpha\in\op{Br}X$, let $X(\AA_k)^\alpha$ be the collection of $x\in X(\AA_k)$ which maps to zero along the composition
	\[X(\AA_k)\stackrel{\op{ev}_\alpha}\to\bigoplus_v\op{Br}k_v\to\QQ/\ZZ.\]
	We then define the \textit{Brauer set} $X(\AA_k)^{\mathrm{Br}}\coloneqq\bigcap_{\alpha\in\op{Br}X}X(\AA_k)^\alpha$. If $X(\AA_k)\ne\emp$ while $X(\AA_k)^{\mathrm{Br}}=\emp$, then we may say that there is a Brauer--Manin obstruction to the Hasse principle.
\end{definition}
\begin{remark}
	Explicitly, $x\in X(\AA_k)$ lives in $X(\AA_k)^\alpha$ if and only if
	\[\sum_v\op{inv}_v\alpha(x_v)=0,\]
	where $\op{inv}_v\colon\op{Br}k_v\to\QQ/\ZZ$ is the local invariant map.
\end{remark}
\begin{remark}
	The preceding discussion shows that $X(k)\subseteq X(\AA_k)^{\mathrm{Br}}$. Thus, $X(\AA_k)^{\mathrm{Br}}=\emp$ tells us that $X(k)=\emp$, so the Brauer--Manin obstruction possibly refines the Hasse principle.
\end{remark}
Next time, we will see an example where the Brauer--Manin obstruction sees $X(\AA_k)^{\mathrm{Br}}=\emp$ even though $X(\AA_k)=\emp$.

\end{document}